% Build:  make report      (or: cd report && pdflatex report.tex, twice)
% Screenshots go in report/screenshots/ with the file names shown in each
% placeholder box (.png or .jpg). They are picked up automatically.
\documentclass[11pt]{article}
\usepackage[margin=1in]{geometry}
\usepackage{graphicx}
\usepackage{float}
\usepackage{booktabs}
\usepackage{enumitem}
\usepackage[hidelinks]{hyperref}

% ---- Fill these in ---------------------------------------------------------
\newcommand{\studentname}{Santiago Soto}
\newcommand{\coursename}{CS 375 -- Operating Systems}
\newcommand{\repolink}{https://github.com/stjamessoto/Module-4-Assignment-3-Lab-Scheduling-Project}
% ----------------------------------------------------------------------------

% \screenshot{file name without extension}{caption}
% Shows screenshots/<name>.png (or .jpg/.jpeg); if it is not there yet, shows a box saying which file to add.
\newcommand{\screenshot}[2]{%
  \begin{figure}[H]
    \centering
    \IfFileExists{screenshots/#1.png}{\includegraphics[width=\linewidth,height=0.6\textheight,keepaspectratio]{screenshots/#1.png}}{%
    \IfFileExists{screenshots/#1.jpg}{\includegraphics[width=\linewidth,height=0.6\textheight,keepaspectratio]{screenshots/#1.jpg}}{%
    \IfFileExists{screenshots/#1.jpeg}{\includegraphics[width=\linewidth,height=0.6\textheight,keepaspectratio]{screenshots/#1.jpeg}}{%
      \fbox{\parbox{0.9\linewidth}{\centering\vspace{1.2cm}
        \textbf{Screenshot not added yet}\\[4pt]
        Save it as \texttt{report/screenshots/\detokenize{#1}.png}\vspace{1.2cm}}}}}}
    \caption{#2}
  \end{figure}}

\title{Lab 8: Task Scheduling Project\\[4pt]\large Lab Exercises and Modular Scheduling Simulator}
\author{\studentname\\\coursename}
\date{\today}

\begin{document}
\maketitle

\noindent\textbf{GitHub repository:}\\\url{\repolink}

\tableofcontents
\newpage

% ============================================================================
\section{Overview}

This report covers both parts of the lab. Part~2 is ten standalone C++ programs, one per
scheduling algorithm, each run on the process table from the assignment
(Table~\ref{tab:processes}). Part~3 combines the same ten algorithms into one command-line
simulator built around a polymorphic \texttt{Scheduler} base class, and uses it to compare
the schedulers on larger workloads.

\begin{table}[H]
  \centering
  \begin{tabular}{lccc}
    \toprule
    Process & Arrival Time & Burst Time & Priority (lower = higher) \\
    \midrule
    P1 & 0 & 8 & 2 \\
    P2 & 1 & 4 & 1 \\
    P3 & 2 & 9 & 3 \\
    P4 & 3 & 5 & 4 \\
    \bottomrule
  \end{tabular}
  \caption{Process table used for every lab exercise.}
  \label{tab:processes}
\end{table}

Every program prints the process table, a text Gantt chart in the form \texttt{ID(duration)}
with a timeline underneath, the waiting and turnaround time of each process, and the
average waiting time, average turnaround time and CPU utilization. CPU utilization is
computed as total burst time divided by elapsed time, so idle time is counted.

% ============================================================================
\section{Part 2: Lab Exercises}

Table~\ref{tab:lab-results} summarises the results of the ten programs on the assignment
table. All ten finish at $t = 26$ with 100\% CPU utilization, because the CPU is never idle.

\begin{table}[H]
  \centering
  \begin{tabular}{llcc}
    \toprule
    Ex. & Scheduler & Avg waiting & Avg turnaround \\
    \midrule
    1 & FCFS & 8.75 & 15.25 \\
    2 & SJF (non-preemptive) & 7.75 & 14.25 \\
    3 & SRTF (preemptive) & 6.50 & 13.00 \\
    4 & Priority with aging & 8.75 & 15.25 \\
    5 & Round Robin ($q = 4$) & 11.75 & 18.25 \\
    6 & Multilevel Queue & 8.75 & 15.25 \\
    7 & Multilevel Feedback Queue & 14.00 & 20.50 \\
    8 & Lottery & \multicolumn{2}{c}{varies with the random seed} \\
    9 & CFS (simplified) & 12.25 & 18.75 \\
    10 & EDF & 6.50 & 13.00 \\
    \bottomrule
  \end{tabular}
  \caption{Lab results on the assignment process table.}
  \label{tab:lab-results}
\end{table}

\subsection{Exercise 1: FCFS (\texttt{fcfs.cpp})}
Processes are sorted by arrival time and each runs to completion. If the next process has
not arrived yet, the CPU idles until it does. The order is P1, P2, P3, P4.
\screenshot{ex01_fcfs}{FCFS output.}

\subsection{Exercise 2: SJF, non-preemptive (\texttt{sjf.cpp})}
Arrived processes are added to a ready queue, the queue is sorted by burst time, and the
shortest one runs to completion. Only P1 is present at $t = 0$, so it runs first; after
that the order is P2, P4, P3. Running the short jobs before P3 lowers the average waiting
time from 8.75 to 7.75.
\screenshot{ex02_sjf}{SJF output.}

\subsection{Exercise 3: SRTF, preemptive (\texttt{srtf.cpp})}
The simulation advances one time unit at a time and always runs the arrived process with
the least remaining time. P2 preempts P1 at $t = 1$ because its remaining time (4) is
shorter than P1's (7). The program logs each preemption, and the Gantt chart starts a new
entry at every context switch. SRTF gives the lowest average waiting time of the ten (6.50).
\screenshot{ex03_srtf}{SRTF output with the preemption logged.}

\subsection{Exercise 4: Priority with aging (\texttt{priority.cpp})}
The ready process with the lowest priority number runs to completion.
\textbf{Enhancement (aging):} for every 5 time units a process has waited, its priority
improves by one level, so a low-priority process cannot starve. The log shows each
process's original and aged priority when it is dispatched; by $t = 21$, P4 has aged from
priority 4 to 1.
\screenshot{ex04_priority}{Priority scheduling output with the aging log.}

\subsection{Exercise 5: Round Robin (\texttt{rr.cpp})}
Processes take turns from a FIFO queue with a quantum of 4. A process that still has work
left is preempted and moved to the back of the queue; processes that arrived during its
time slice are queued ahead of it. Round Robin has a higher average waiting time (11.75)
than FCFS here, but every process gets on the CPU much sooner.
\screenshot{ex05_rr}{Round Robin output.}

\subsection{Exercise 6: Multilevel Queue (\texttt{mlq.cpp})}
Processes with priority below 3 (P1, P2) go to the high queue, which uses Round Robin with
a quantum of 4; the rest (P3, P4) go to the low queue, which is FCFS. The high queue is
always served first, and the low queue runs one time unit at a time so a high-priority
arrival preempts it immediately.
\screenshot{ex06_mlq}{Multilevel Queue output showing the queue assignments.}

\subsection{Exercise 7: Multilevel Feedback Queue (\texttt{mlfq.cpp})}
There are three Round Robin queues with quanta 2, 4 and 8, and every process starts in the
top queue. \textbf{Feedback:} a process that uses its whole quantum is treated as CPU-bound
and demoted one level. \textbf{Enhancement:} a process that has waited 8 time units in a
lower queue is promoted one level. The feedback log lists every demotion and promotion.
\screenshot{ex07_mlfq}{MLFQ output with the feedback log.}

\subsection{Exercise 8: Lottery (\texttt{lottery.cpp})}
Each process gets $10 / \mathit{priority}$ tickets (P1 = 5, P2 = 10, P3 = 3, P4 = 2). At
each step a random ticket is drawn among the ready processes and the winner runs for a
quantum of 2. \textbf{Enhancement (dynamic tickets):} every ready process that loses a draw
gains one ticket, and the winner returns to its base amount. The program accepts an
optional seed (\texttt{./lottery 7}) so that a run can be repeated exactly.
\screenshot{ex08_lottery}{Lottery output showing every draw and the ticket counts.}

\subsection{Exercise 9: Simplified CFS (\texttt{cfs.cpp})}
Ready processes are kept in a min-heap ordered by virtual runtime, and the process with
the smallest value runs for a slice of 2. \textbf{Enhancement (weights):} each process has
a weight of $1024 / 1.25^{\,\mathit{priority}-1}$ and its virtual runtime grows by
$\mathit{slice} \times 1024 / \mathit{weight}$, so a higher-priority process accumulates
virtual runtime more slowly and receives more CPU time. New arrivals start at the current
minimum virtual runtime.
\screenshot{ex09_cfs}{CFS output with weights and final virtual runtimes.}

\subsection{Exercise 10: EDF (\texttt{edf.cpp})}
Each process gets a deadline of $\mathit{arrival} + 2 \times \mathit{burst}$ (P1 = 16,
P2 = 9, P3 = 20, P4 = 13). A priority queue ordered by deadline picks the process to run
each time unit, so a nearer deadline preempts. \textbf{Enhancement (schedulability):} the
program first computes the density $\sum \mathit{burst} / (\mathit{deadline} -
\mathit{arrival})$, which is 2.0 here, so the deadlines are not guaranteed; after the run it
reports each deadline as met or missed. P1 and P3 miss theirs.
\screenshot{ex10_edf}{EDF output with the schedulability check and deadline report.}

% ============================================================================
\section{Part 3: Modular Task Scheduling Simulator}

\subsection{Design}
The simulator is built around an abstract base class with one pure virtual method:

\begin{verbatim}
class Scheduler {
public:
    virtual ~Scheduler() = default;
    virtual void schedule(std::vector<Process>& processes,
                          Gantt& gantt, int& total_time) = 0;
    virtual std::string name() const = 0;
};
\end{verbatim}

Each of the ten algorithms is a derived class that overrides \texttt{schedule()}.
\texttt{main()} selects one by name through a small factory function and only ever uses it
through a \texttt{std::unique\_ptr<Scheduler>}. Adding a scheduler means adding one class
and one line in the factory. The base class also provides the helpers every scheduler
needs: running a process for a time slice, recording Gantt entries, and inserting idle time.

\begin{table}[H]
  \centering
  \begin{tabular}{ll}
    \toprule
    File & Contents \\
    \midrule
    \texttt{simulator.cpp} & \texttt{main()}, argument parsing, scheduler factory, reports \\
    \texttt{process.h} & the \texttt{Process} struct \\
    \texttt{scheduler.h} & the abstract \texttt{Scheduler} base class and shared helpers \\
    \texttt{schedulers.h} & the ten derived scheduler classes \\
    \texttt{utils.h} & input file loading, task generation, metrics, printing \\
    \texttt{tasks.txt}, \texttt{tasks\_deadlines.txt} & sample input files \\
    \bottomrule
  \end{tabular}
  \caption{Source files of the simulator.}
\end{table}

\subsection{Usage}
\begin{verbatim}
g++ -std=c++17 -pthread simulator.cpp -o simulator
./simulator --scheduler rr --quantum 4
./simulator --scheduler srtf --input tasks.txt --output log.txt
./simulator --compare --random --num 20 --seed 42
\end{verbatim}

\begin{table}[H]
  \centering
  \begin{tabular}{ll}
    \toprule
    Option & Meaning \\
    \midrule
    \texttt{--scheduler <name>} & fcfs, sjf, srtf, priority, rr, mlq, mlfq, lottery, cfs, edf \\
    \texttt{--compare [a,b,...]} & run several (default: all) schedulers on the same tasks \\
    \texttt{--input <file>} & tasks from a file: \texttt{ID arrival burst priority [deadline]} \\
    \texttt{--random}, \texttt{--num <n>} & generate $n$ random tasks (default 10) \\
    \texttt{--web} & generate a web-server request workload \\
    \texttt{--seed <n>} & seed for random tasks, lottery draws and I/O \\
    \texttt{--quantum <n>} & quantum for rr, mlq and lottery (default 4) \\
    \texttt{--aging <n>} & priority aging interval (default 5, 0 turns it off) \\
    \texttt{--io <percent>} & rr: chance of blocking for I/O after a time slice \\
    \texttt{--cpus <n>} & number of CPUs, each simulated in its own thread \\
    \texttt{--output <file>} & also write the results to a file \\
    \bottomrule
  \end{tabular}
  \caption{Command-line options.}
\end{table}

With no task option the simulator uses the assignment table. Besides the Gantt chart it
reports average waiting, turnaround and response time, CPU utilization, throughput and
the number of context switches.

\subsection{Enhancements}
\begin{itemize}[leftmargin=*]
  \item \textbf{Priority aging.} A waiting process gains one priority level every
        \texttt{--aging} time units. The scheduler logs each decision that aging changed.
  \item \textbf{EDF deadlines.} Deadlines are read from the optional fifth column of the
        input file; tasks without one get $\mathit{arrival} + 2 \times \mathit{burst}$. The
        scheduler reports a density-based schedulability estimate and every missed deadline.
  \item \textbf{I/O simulation (Round Robin).} With \texttt{--io p}, a process that has not
        finished at the end of its time slice blocks for 1--5 time units with probability
        $p$\%, then rejoins the ready queue. Time spent blocked is not counted as waiting.
  \item \textbf{MLFQ feedback, lottery dynamic tickets, CFS weights.} The same enhancements
        as in the lab exercises.
  \item \textbf{Web-server scenario and load balancing.} \texttt{--web} generates a stream
        of requests in which about one in five is long. With \texttt{--cpus n}, a load
        balancer sends each request to the CPU that will be free soonest, and each CPU runs
        its own scheduler instance in a separate \texttt{std::thread}.
  \item \textbf{Error handling.} Unknown options and schedulers, missing values,
        non-numeric or out-of-range numbers and unreadable files produce an error message
        and a non-zero exit code. Malformed lines in an input file are reported with their
        line number and skipped.
\end{itemize}

\subsection{Test Runs}

\screenshot{p3_01_rr_default}{Round Robin on the default tasks
(\texttt{--scheduler rr --quantum 4}). The result matches the lab program.}

\screenshot{p3_02_rr_random}{Round Robin on ten random tasks
(\texttt{--scheduler rr --random --seed 42}).}

\screenshot{p3_03_srtf_file}{SRTF on tasks read from \texttt{tasks.txt}.}

\screenshot{p3_04_compare}{All ten schedulers compared on \texttt{tasks.txt}.}

\screenshot{p3_05_priority_aging}{Priority scheduling with aging every 3 units; the events
list shows where aging changed the order.}

\screenshot{p3_06_edf_deadlines}{EDF with deadlines read from \texttt{tasks\_deadlines.txt}.}

\screenshot{p3_07_rr_io}{Round Robin with I/O simulation (\texttt{--io 40}).}

\screenshot{p3_08_web_cpus}{Web-server workload load-balanced across two CPUs.}

\screenshot{p3_09_output_file}{Writing the results to a file with \texttt{--output}.}

\screenshot{p3_10_errors}{Error handling for an unknown scheduler, a missing file and an
invalid number.}

\subsection{Comparison and Analysis}

Table~\ref{tab:compare} shows all ten schedulers on the 12-task sample workload in
\texttt{tasks.txt}, which mixes long jobs (bursts of 9--12) with short ones (bursts of
1--3). It was produced with
\texttt{./simulator --compare --input tasks.txt --seed 42}. CPU utilization (100\%) and
throughput (0.2034 tasks per time unit) are the same for every scheduler, because on one
CPU with no idle time the order of the work does not change when all of it finishes.

\begin{table}[H]
  \centering
  \begin{tabular}{lrrrr}
    \toprule
    Scheduler & Avg waiting & Avg turnaround & Avg response & Context switches \\
    \midrule
    fcfs & 22.67 & 27.58 & 22.67 & 11 \\
    sjf & 12.50 & 17.42 & 12.50 & 11 \\
    srtf & 9.92 & 14.83 & 7.08 & 14 \\
    priority (aging 5) & 19.83 & 24.75 & 19.83 & 11 \\
    rr ($q = 4$) & 22.17 & 27.08 & 12.08 & 19 \\
    mlq & 13.00 & 17.92 & 12.08 & 15 \\
    mlfq & 20.25 & 25.17 & 3.25 & 29 \\
    lottery (seed 42) & 21.08 & 26.00 & 11.67 & 19 \\
    cfs & 18.17 & 23.08 & 2.25 & 30 \\
    edf & 10.75 & 15.67 & 8.75 & 14 \\
    \bottomrule
  \end{tabular}
  \caption{All schedulers on \texttt{tasks.txt} (12 tasks, one CPU).}
  \label{tab:compare}
\end{table}

\paragraph{Round Robin vs.\ SRTF.} SRTF has less than half the average waiting time of
Round Robin (9.92 against 22.17) and a lower average response time (7.08 against 12.08),
with fewer context switches (14 against 19). SRTF wins on these averages because it always
finishes the shortest work first. Its weaknesses do not show in the averages: it needs to
know each burst length in advance, which a real system can only estimate, and a long job
can be postponed indefinitely while shorter ones keep arriving. Round Robin needs no
burst estimates and bounds how long any task waits for its next turn, which is why it suits
interactive systems, at the cost of more context switches and longer turnaround.

\paragraph{Response time vs.\ waiting time.} MLFQ and CFS have the lowest average response
times (3.25 and 2.25) because every new task gets a short slice almost immediately, but
they have the highest context-switch counts (29 and 30) and much higher average waiting
times than SJF or SRTF, since long jobs are interrupted repeatedly. FCFS is the opposite: the fewest
switches, and the worst waiting and response time (22.67), because short tasks queue behind
long ones (the convoy effect).

\paragraph{Effect of aging.} On the same workload, priority scheduling without aging
(\texttt{--aging 0}) has an average waiting time of 14.58, but the low-priority task T5
waits 38 time units. With aging every 5 units T5 waits 20, and with aging every 3 units it
waits 15, while the average rises to 19.83 and 22.67. Aging trades average waiting time for
a bound on how long a low-priority task is postponed.

\paragraph{CPU-bound vs.\ I/O-bound workloads.} For CPU-bound batch work, where only total
completion time matters and nobody is waiting on a response, the non-preemptive schedulers
are the best fit: SJF gives a low average turnaround (17.42) with the minimum number of
context switches, and FCFS is acceptable when the jobs are of similar length. For
I/O-bound and interactive work, tasks use the CPU in short bursts and need it again quickly
after each I/O operation, so response time is the metric that matters. MLFQ and CFS are the
best fit there: a task that gives up the CPU before its quantum ends stays in a
high-priority queue (MLFQ) or keeps a low virtual runtime (CFS), so it is scheduled soon
after it becomes ready, and CPU-bound tasks sink to the lower queues without being starved.
Round Robin is a reasonable simple choice for interactive work but treats both kinds of
task the same. In the Round Robin I/O simulation on this workload (\texttt{--io 40
--seed 7}), two tasks blocked for I/O and the other tasks used the CPU in the meantime, so
utilization stayed at 100\%. EDF is the right choice only when tasks have real deadlines: on
\texttt{tasks\_deadlines.txt} it meets all 10.

\paragraph{Load balancing.} For the 12-request web workload (\texttt{--web --num 12
--seed 3}) under Round Robin, one CPU finishes at $t = 51$ with an average waiting time of
13.25. Two CPUs finish at $t = 28$ with an average waiting time of 4.00, at 89.29\% CPU
utilization.

\section{References}
\begin{enumerate}[leftmargin=*]
  \item Lab handout, \emph{Lab 8: Implementing and Enhancing Task Schedulers in C++} ---
        source of the starter code (\texttt{Process} struct, \texttt{Scheduler} base class,
        Round Robin example, argument parsing).
  \item A.~Silberschatz, P.~B.~Galvin, G.~Gagne, \emph{Operating System Concepts}, 10th
        ed., Wiley, 2018. Chapter 5, CPU Scheduling.
  \item R.~H.~Arpaci-Dusseau, A.~C.~Arpaci-Dusseau, \emph{Operating Systems: Three Easy
        Pieces}, chapters 7--9 (scheduling, MLFQ, lottery scheduling and CFS).
        \url{https://pages.cs.wisc.edu/~remzi/OSTEP/}
  \item C.~A.~Waldspurger, W.~E.~Weihl, ``Lottery Scheduling: Flexible Proportional-Share
        Resource Management,'' \emph{OSDI}, 1994.
  \item C.~L.~Liu, J.~W.~Layland, ``Scheduling Algorithms for Multiprogramming in a
        Hard-Real-Time Environment,'' \emph{Journal of the ACM}, 20(1), 1973.
\end{enumerate}

\end{document}
